% =============================================================================
% Clientes e usuarios do projeto (projeto Riachuelo, PRO3474).
%
% Consolida a definicao dos clientes do projeto e a segmentacao dos
% consumidores investigados. O mercado e o segmento-alvo B2B estao na secao
% anterior (definicao_mercado.tex); ciclo de vida tera secao propria.
% Os arquivos 5-Refinar_ciclo_de_vida.tex e 6-segmentacao_quatro_grupos.tex
% permanecem na pasta, fora do documento.tex, para reaproveitamento.
%
% Labels preservados (referenciados por outras secoes):
%   subsec:segmentacao_clientes_usuarios, subsec:clientes_lacunas,
%   subsubsec:segmentacao_hipoteses
% =============================================================================

\section{CLIENTES E USUÁRIOS DO PROJETO}
\label{sec:clientes_usuarios}

\subsection{Cliente, usuário e outros atores}
\label{subsec:conceitos_cliente_usuario}

O projeto deverá operar em dois níveis, que a análise precisa manter separados. No primeiro, a solução de trava antifurto conectada ao aplicativo será vendida em modelo B2B à Riachuelo, a organização que decidirá adquirir e implantar o sistema. No segundo, dentro da loja, a mesma solução servirá à jornada de compra do consumidor, que escolhe, paga e leva a peça de roupa. Por esse motivo, cada análise deste relatório indica a qual dos dois níveis se refere.

Para distinguir os atores envolvidos, a equipe adotou a classificação de \citeonline{rozenfeld2006gestao}. No nível da jornada de compra, os clientes externos são os consumidores que compram roupa na loja, externos à Riachuelo e incluídos no escopo de uso da solução. Os clientes internos são as áreas da Riachuelo que operam o sistema ou decidem sobre ele sem comprar roupa na loja. O termo usuário, por sua vez, designa quem manuseia o produto físico e o aplicativo no dia a dia. Neste projeto, tanto o consumidor quanto parte do cliente interno serão usuários, cada um em um ponto diferente da jornada, o que distingue o caso de um produto de consumo típico, em que cliente e usuário costumam ser a mesma pessoa.

A mesma classificação separa comprador e consumidor. Em geral, quem paga pela peça também a veste, mas um adulto pode comprar roupa infantil para uma criança e, nesse caso, quem interage com o aplicativo não é quem usa a peça. Como essa distinção ainda não havia sido investigada em campo, ela ficou registrada como pergunta em aberto (Seção~\ref{subsec:clientes_lacunas}). Por fim, Rozenfeld et al. também preveem o cliente intermediário, papel típico dos canais de distribuição entre fabricante e consumidor final. A equipe não identificou ator equivalente neste projeto, já que a Riachuelo opera a própria loja onde a solução funcionará, e por isso a categoria não foi aplicada.

\subsection{Cliente externo: o varejista}
\label{subsec:cliente_externo}

No nível B2B, o cliente externo corresponde às empresas do varejo de roupas que adotarem o sistema, pois serão elas que decidirão a aquisição e pagarão pela solução. A Riachuelo serve de referência para o desenvolvimento, e outros varejistas com necessidades semelhantes poderão adotar o mesmo sistema.

\subsection{Usuários externos: os consumidores da loja}
\label{subsec:usuarios_externos}

Os consumidores deverão usar o aplicativo para identificar a peça pelo código associado à etiqueta, realizar ou confirmar o pagamento e retirar a peça após a confirmação. Serão, portanto, usuários e clientes externos ao mesmo tempo, já que decidem a compra e manuseiam o produto em seguida.

\subsection{Clientes internos: operação e gestão da Riachuelo}
\label{subsec:clientes_internos}

Duas frentes da Riachuelo deverão interagir com a solução sem comprar roupa na loja. A operação reúne vendas, estoque, reposição, caixa e segurança, funções que tocam a loja no dia a dia e que manuseiam, monitoram ou sofrem o impacto direto de qualquer mudança no processo físico, efeito que não aparece na voz do consumidor. A gestão, por sua vez, reúne a média e a alta gerência, que avaliarão a solução por indicadores, resultado operacional, retorno econômico, risco e impacto estratégico. Em um modelo B2B, é esse grupo que aprova a compra da solução, decisão da qual quem usa o produto na loja não participa.

Nesta etapa, a equipe não investigou nenhuma das duas frentes, porque ainda não tinha acesso estruturado à operação interna para conduzir entrevistas ou surveys consistentes com esses grupos. Ambas permaneceram registradas como atores relevantes, a investigar quando esse acesso existir.

\subsection{Consumidores investigados: Cliente Premium e Cliente Popular}
\label{subsec:segmentacao_clientes_usuarios}

O esforço de campo desta etapa concentrou-se em dois grupos de consumidores, definidos pelo formato da loja frequentada. O Cliente Premium corresponde ao consumidor de lojas ou formatos voltados a público de maior poder aquisitivo e interage com a jornada completa em estudo: escolhe o produto, paga e, caso a solução avance, sairá da loja sem passar pelo caixa. O Cliente Popular corresponde ao consumidor do varejo de massa, também usuário final da jornada. A observação de campo apontou, entre os dois, expectativas diferentes de autonomia, atendimento e tecnologia, com um padrão de compra mais assistido pelo vendedor no Cliente Popular, inclusive na oferta de crédito. O que já havia sido observado em campo e o que ainda era hipótese a testar está resumido na Tabela~\ref{tab:segmentacao_grupos}.

\begin{table}[htbp]
    \centering
    \caption{Cliente Premium e Cliente Popular: observações e hipóteses}
    \label{tab:segmentacao_grupos}
    \footnotesize
    \renewcommand{\arraystretch}{1.4}
    \begin{tabularx}{\textwidth}{@{}
        >{\RaggedRight\arraybackslash}p{3.1cm}
        >{\RaggedRight\arraybackslash}X
        >{\RaggedRight\arraybackslash}X
        >{\centering\arraybackslash}p{1.9cm}@{}}
        \toprule
        \textbf{Dimensão} & \textbf{Cliente Premium} & \textbf{Cliente Popular} & \textbf{Base} \\
        \midrule
        Papel do vendedor
            & Menor influência sobre a decisão de cartão de crédito
            & Maior influência sobre a decisão de cartão de crédito
            & Campo \\
        \addlinespace
        Autoatendimento
            & Maior preferência declarada
            & Menor preferência declarada
            & Campo \\
        \addlinespace
        Dependência do vendedor
            & Pode estar diminuindo, com valorização crescente de autonomia
            & Padrão ainda não medido, pode seguir tendência semelhante em outro ritmo
            & Hipótese \\
        \addlinespace
        Lógica de exposição de produto
            & Não caracterizada nesta etapa
            & Prefere ver muitas peças expostas (lógica de abundância)
            & Campo \\
        \bottomrule
    \end{tabularx}
    \source{Elaborado pelos autores (2026).}
\end{table}

Já nessa observação inicial, a preferência entre autonomia e atendimento assistido pareceu variar conforme o momento da jornada e o grupo de consumidor. Por isso, a equipe evitou tratar a autonomia como vantagem automática, o que poderia levar a uma solução que retira o suporte no momento em que o cliente ainda precisa dele.

\subsection{Hipóteses de pesquisa}
\label{subsec:clientes_lacunas}
\label{subsubsec:segmentacao_hipoteses}

As hipóteses a seguir orientaram a pesquisa com Cliente Premium e Cliente Popular. Elas serviram como pontos de partida da investigação, sujeitos a confirmação, refutação ou ajuste, e a definição do problema do projeto ficou reservada a uma etapa posterior.

\begin{itemize}
    \item H1: no Cliente Premium, a dependência do vendedor para a decisão de compra está diminuindo, com valorização crescente de autonomia durante a jornada.
    \item H2: o mesmo movimento pode ocorrer no Cliente Popular, em velocidade ou intensidade diferente.
    \item H3: a fricção mais relevante da jornada de compra pode incluir disponibilidade de produto, dificuldade de encontrar itens na loja ou integração entre canal físico e digital, além de fila e pagamento.
    \item H4: os dois grupos valorizam atendimento humano em momentos distintos da jornada, como a escolha e a prova de roupa em contraste com o pagamento, e essa valorização muda ao longo do percurso de compra.
\end{itemize}

Além das hipóteses, dois pontos permaneceram em aberto por falta de dado de campo. O primeiro era saber se comprador e consumidor coincidem na maioria dos casos investigados ou se situações como a compra de roupa infantil por um adulto aparecem com frequência relevante. O segundo era saber se a segmentação entre Cliente Premium e Cliente Popular captura as diferenças que mais importam para o produto ou se outras dimensões, como a familiaridade com tecnologia, explicam melhor o comportamento observado. Os resultados referentes às hipóteses e a esses dois pontos estão discutidos na Seção~\ref{sec:sintese}.
